\documentclass[10pt]{article}
% Hojas resumen: un A4 por delante y por detrás, dos columnas por cara.
% Uso: \documentclass[10pt]{article} % Hojas resumen: un A4 por delante y por detrás, dos columnas por cara.
% Uso: \documentclass[10pt]{article} % Hojas resumen: un A4 por delante y por detrás, dos columnas por cara.
% Uso: \documentclass[10pt]{article} \input{estilo-resumen.tex} \bloque{analisis|algebra|azar} \cabecera{n}{Título}{Lecciones}
\usepackage[utf8]{inputenc}
\usepackage[T1]{fontenc}
\usepackage[spanish,es-nodecimaldot]{babel}
\usepackage{lmodern}
\usepackage{amsmath,amssymb}
\usepackage{graphicx}
\usepackage[a4paper,left=10mm,right=10mm,top=15mm,bottom=13mm,headheight=11pt,headsep=3mm,footskip=6mm]{geometry}
\usepackage{multicol}
\setlength{\columnsep}{6mm}
\setlength{\columnseprule}{0.3pt}
\usepackage{xcolor}
\usepackage{enumitem}
\setlist{nosep,leftmargin=1.1em,itemsep=1pt}
\usepackage{array}
\usepackage{needspace}
\usepackage{pgfplots}
\usepgfplotslibrary{fillbetween}
\pgfplotsset{compat=1.18,
  mini/.style={width=4.3cm,height=3.3cm,scale only axis=false,axis lines=middle,xtick=\empty,ytick=\empty,samples=120,
    xlabel={\scriptsize$x$},ylabel={\scriptsize$y$},title style={font=\scriptsize,yshift=-4pt},clip=true,
    every axis plot/.append style={thick,color=acento}}}
\usepackage{ragged2e}
\usepackage[most]{tcolorbox}
\usepackage{fancyhdr,lastpage}
\usepackage[hidelinks,pdfencoding=auto]{hyperref}
\hypersetup{pdflang=es-ES,pdfsubject={Matemáticas Aplicadas a las Ciencias Sociales II}}

\definecolor{analisis}{HTML}{1D4ED8}
\definecolor{algebra}{HTML}{86198F}
\definecolor{azar}{HTML}{9A3412}
\definecolor{aviso}{HTML}{B91C1C}
\colorlet{acento}{analisis}
\newcommand{\bloque}[1]{\colorlet{acento}{#1}}

\setlength{\parindent}{0pt}
\setlength{\parskip}{2pt}
\renewcommand{\arraystretch}{1.2}
\AtBeginDocument{\RaggedRight\fontsize{9.8}{11.9}\selectfont\setlength{\abovedisplayskip}{3pt}\setlength{\belowdisplayskip}{3pt}\setlength{\abovedisplayshortskip}{1pt}\setlength{\belowdisplayshortskip}{1pt}}

\newcommand{\temanum}{}\newcommand{\tematitulo}{}
\pagestyle{fancy}
\fancyhf{}
\fancyhead[L]{\footnotesize\color{acento}\textbf{Hoja resumen}\quad Tema \temanum. \tematitulo}
\fancyhead[R]{\footnotesize Matemáticas Aplicadas a las Ciencias Sociales II --- 2.\textsuperscript{o} Bachillerato}
\fancyfoot[L]{\scriptsize IES Al-Ándalus --- Juan María Gámez Ortiz, con la ayuda de Claude y de Codex --- CC BY-NC 4.0}
\fancyfoot[R]{\scriptsize Cara \thepage\ de \pageref{LastPage}}
\renewcommand{\headrulewidth}{0.4pt}

% Cabecera de la primera cara: número, título y lecciones del tema.
\newcommand{\cabecera}[3]{%
  \renewcommand{\temanum}{#1}\renewcommand{\tematitulo}{#2}%
  \hypersetup{pdftitle={Hoja resumen. Tema #1. #2}}%
  \noindent\begin{tcolorbox}[enhanced,colback=acento!8,colframe=acento,boxrule=0.6pt,arc=2pt,left=4pt,right=4pt,top=2pt,bottom=2pt]
  \begin{minipage}[c]{1.1cm}\includegraphics[width=1cm]{logo-al-andalus.png}\end{minipage}%
  \begin{minipage}[c]{\dimexpr\linewidth-1.2cm}{\large\bfseries\color{acento}Tema #1. #2}\\[1pt]{\scriptsize #3}\end{minipage}
  \end{tcolorbox}\vspace{-2pt}}

% Apartado con filete del color del bloque.
\newcommand{\seccion}[1]{\par\needspace{5\baselineskip}\vspace{3pt}{\color{acento}\bfseries\normalsize #1}\par\vspace{-1pt}{\color{acento}\hrule height 0.5pt}\vspace{2pt}}
\newcommand{\sub}[1]{\par\needspace{3\baselineskip}\vspace{1pt}\textbf{\color{acento}#1}\enspace}

% Fórmula clave, aviso de errores y ejemplo.
\tcbset{before upper=\RaggedRight}
\newtcolorbox{formula}{enhanced,colback=acento!6,colframe=acento!70,boxrule=0.5pt,arc=1.5pt,left=3pt,right=3pt,top=0pt,bottom=0pt,before skip=2pt,after skip=2pt}
\newtcolorbox{ojo}{enhanced,colback=aviso!5,colframe=aviso!60,boxrule=0.5pt,arc=1.5pt,left=3pt,right=3pt,top=1pt,bottom=1pt,before skip=3pt,after skip=2pt,
  fonttitle=\bfseries\footnotesize,coltitle=aviso,colbacktitle=aviso!5,title={Ojo},attach title to upper={\enspace}}
\newtcolorbox{ejemplo}{enhanced,colback=white,colframe=acento!45,boxrule=0.4pt,arc=1.5pt,left=3pt,right=3pt,top=1pt,bottom=1pt,before skip=3pt,after skip=2pt,
  fonttitle=\bfseries\footnotesize,coltitle=acento,colbacktitle=white,title={Ejemplo},attach title to upper={\enspace}}
\newcommand{\dx}{\,dx}
 \bloque{analisis|algebra|azar} \cabecera{n}{Título}{Lecciones}
\usepackage[utf8]{inputenc}
\usepackage[T1]{fontenc}
\usepackage[spanish,es-nodecimaldot]{babel}
\usepackage{lmodern}
\usepackage{amsmath,amssymb}
\usepackage{graphicx}
\usepackage[a4paper,left=10mm,right=10mm,top=15mm,bottom=13mm,headheight=11pt,headsep=3mm,footskip=6mm]{geometry}
\usepackage{multicol}
\setlength{\columnsep}{6mm}
\setlength{\columnseprule}{0.3pt}
\usepackage{xcolor}
\usepackage{enumitem}
\setlist{nosep,leftmargin=1.1em,itemsep=1pt}
\usepackage{array}
\usepackage{needspace}
\usepackage{pgfplots}
\usepgfplotslibrary{fillbetween}
\pgfplotsset{compat=1.18,
  mini/.style={width=4.3cm,height=3.3cm,scale only axis=false,axis lines=middle,xtick=\empty,ytick=\empty,samples=120,
    xlabel={\scriptsize$x$},ylabel={\scriptsize$y$},title style={font=\scriptsize,yshift=-4pt},clip=true,
    every axis plot/.append style={thick,color=acento}}}
\usepackage{ragged2e}
\usepackage[most]{tcolorbox}
\usepackage{fancyhdr,lastpage}
\usepackage[hidelinks,pdfencoding=auto]{hyperref}
\hypersetup{pdflang=es-ES,pdfsubject={Matemáticas Aplicadas a las Ciencias Sociales II}}

\definecolor{analisis}{HTML}{1D4ED8}
\definecolor{algebra}{HTML}{86198F}
\definecolor{azar}{HTML}{9A3412}
\definecolor{aviso}{HTML}{B91C1C}
\colorlet{acento}{analisis}
\newcommand{\bloque}[1]{\colorlet{acento}{#1}}

\setlength{\parindent}{0pt}
\setlength{\parskip}{2pt}
\renewcommand{\arraystretch}{1.2}
\AtBeginDocument{\RaggedRight\fontsize{9.8}{11.9}\selectfont\setlength{\abovedisplayskip}{3pt}\setlength{\belowdisplayskip}{3pt}\setlength{\abovedisplayshortskip}{1pt}\setlength{\belowdisplayshortskip}{1pt}}

\newcommand{\temanum}{}\newcommand{\tematitulo}{}
\pagestyle{fancy}
\fancyhf{}
\fancyhead[L]{\footnotesize\color{acento}\textbf{Hoja resumen}\quad Tema \temanum. \tematitulo}
\fancyhead[R]{\footnotesize Matemáticas Aplicadas a las Ciencias Sociales II --- 2.\textsuperscript{o} Bachillerato}
\fancyfoot[L]{\scriptsize IES Al-Ándalus --- Juan María Gámez Ortiz, con la ayuda de Claude y de Codex --- CC BY-NC 4.0}
\fancyfoot[R]{\scriptsize Cara \thepage\ de \pageref{LastPage}}
\renewcommand{\headrulewidth}{0.4pt}

% Cabecera de la primera cara: número, título y lecciones del tema.
\newcommand{\cabecera}[3]{%
  \renewcommand{\temanum}{#1}\renewcommand{\tematitulo}{#2}%
  \hypersetup{pdftitle={Hoja resumen. Tema #1. #2}}%
  \noindent\begin{tcolorbox}[enhanced,colback=acento!8,colframe=acento,boxrule=0.6pt,arc=2pt,left=4pt,right=4pt,top=2pt,bottom=2pt]
  \begin{minipage}[c]{1.1cm}\includegraphics[width=1cm]{logo-al-andalus.png}\end{minipage}%
  \begin{minipage}[c]{\dimexpr\linewidth-1.2cm}{\large\bfseries\color{acento}Tema #1. #2}\\[1pt]{\scriptsize #3}\end{minipage}
  \end{tcolorbox}\vspace{-2pt}}

% Apartado con filete del color del bloque.
\newcommand{\seccion}[1]{\par\needspace{5\baselineskip}\vspace{3pt}{\color{acento}\bfseries\normalsize #1}\par\vspace{-1pt}{\color{acento}\hrule height 0.5pt}\vspace{2pt}}
\newcommand{\sub}[1]{\par\needspace{3\baselineskip}\vspace{1pt}\textbf{\color{acento}#1}\enspace}

% Fórmula clave, aviso de errores y ejemplo.
\tcbset{before upper=\RaggedRight}
\newtcolorbox{formula}{enhanced,colback=acento!6,colframe=acento!70,boxrule=0.5pt,arc=1.5pt,left=3pt,right=3pt,top=0pt,bottom=0pt,before skip=2pt,after skip=2pt}
\newtcolorbox{ojo}{enhanced,colback=aviso!5,colframe=aviso!60,boxrule=0.5pt,arc=1.5pt,left=3pt,right=3pt,top=1pt,bottom=1pt,before skip=3pt,after skip=2pt,
  fonttitle=\bfseries\footnotesize,coltitle=aviso,colbacktitle=aviso!5,title={Ojo},attach title to upper={\enspace}}
\newtcolorbox{ejemplo}{enhanced,colback=white,colframe=acento!45,boxrule=0.4pt,arc=1.5pt,left=3pt,right=3pt,top=1pt,bottom=1pt,before skip=3pt,after skip=2pt,
  fonttitle=\bfseries\footnotesize,coltitle=acento,colbacktitle=white,title={Ejemplo},attach title to upper={\enspace}}
\newcommand{\dx}{\,dx}
 \bloque{analisis|algebra|azar} \cabecera{n}{Título}{Lecciones}
\usepackage[utf8]{inputenc}
\usepackage[T1]{fontenc}
\usepackage[spanish,es-nodecimaldot]{babel}
\usepackage{lmodern}
\usepackage{amsmath,amssymb}
\usepackage{graphicx}
\usepackage[a4paper,left=10mm,right=10mm,top=15mm,bottom=13mm,headheight=11pt,headsep=3mm,footskip=6mm]{geometry}
\usepackage{multicol}
\setlength{\columnsep}{6mm}
\setlength{\columnseprule}{0.3pt}
\usepackage{xcolor}
\usepackage{enumitem}
\setlist{nosep,leftmargin=1.1em,itemsep=1pt}
\usepackage{array}
\usepackage{needspace}
\usepackage{pgfplots}
\usepgfplotslibrary{fillbetween}
\pgfplotsset{compat=1.18,
  mini/.style={width=4.3cm,height=3.3cm,scale only axis=false,axis lines=middle,xtick=\empty,ytick=\empty,samples=120,
    xlabel={\scriptsize$x$},ylabel={\scriptsize$y$},title style={font=\scriptsize,yshift=-4pt},clip=true,
    every axis plot/.append style={thick,color=acento}}}
\usepackage{ragged2e}
\usepackage[most]{tcolorbox}
\usepackage{fancyhdr,lastpage}
\usepackage[hidelinks,pdfencoding=auto]{hyperref}
\hypersetup{pdflang=es-ES,pdfsubject={Matemáticas Aplicadas a las Ciencias Sociales II}}

\definecolor{analisis}{HTML}{1D4ED8}
\definecolor{algebra}{HTML}{86198F}
\definecolor{azar}{HTML}{9A3412}
\definecolor{aviso}{HTML}{B91C1C}
\colorlet{acento}{analisis}
\newcommand{\bloque}[1]{\colorlet{acento}{#1}}

\setlength{\parindent}{0pt}
\setlength{\parskip}{2pt}
\renewcommand{\arraystretch}{1.2}
\AtBeginDocument{\RaggedRight\fontsize{9.8}{11.9}\selectfont\setlength{\abovedisplayskip}{3pt}\setlength{\belowdisplayskip}{3pt}\setlength{\abovedisplayshortskip}{1pt}\setlength{\belowdisplayshortskip}{1pt}}

\newcommand{\temanum}{}\newcommand{\tematitulo}{}
\pagestyle{fancy}
\fancyhf{}
\fancyhead[L]{\footnotesize\color{acento}\textbf{Hoja resumen}\quad Tema \temanum. \tematitulo}
\fancyhead[R]{\footnotesize Matemáticas Aplicadas a las Ciencias Sociales II --- 2.\textsuperscript{o} Bachillerato}
\fancyfoot[L]{\scriptsize IES Al-Ándalus --- Juan María Gámez Ortiz, con la ayuda de Claude y de Codex --- CC BY-NC 4.0}
\fancyfoot[R]{\scriptsize Cara \thepage\ de \pageref{LastPage}}
\renewcommand{\headrulewidth}{0.4pt}

% Cabecera de la primera cara: número, título y lecciones del tema.
\newcommand{\cabecera}[3]{%
  \renewcommand{\temanum}{#1}\renewcommand{\tematitulo}{#2}%
  \hypersetup{pdftitle={Hoja resumen. Tema #1. #2}}%
  \noindent\begin{tcolorbox}[enhanced,colback=acento!8,colframe=acento,boxrule=0.6pt,arc=2pt,left=4pt,right=4pt,top=2pt,bottom=2pt]
  \begin{minipage}[c]{1.1cm}\includegraphics[width=1cm]{logo-al-andalus.png}\end{minipage}%
  \begin{minipage}[c]{\dimexpr\linewidth-1.2cm}{\large\bfseries\color{acento}Tema #1. #2}\\[1pt]{\scriptsize #3}\end{minipage}
  \end{tcolorbox}\vspace{-2pt}}

% Apartado con filete del color del bloque.
\newcommand{\seccion}[1]{\par\needspace{5\baselineskip}\vspace{3pt}{\color{acento}\bfseries\normalsize #1}\par\vspace{-1pt}{\color{acento}\hrule height 0.5pt}\vspace{2pt}}
\newcommand{\sub}[1]{\par\needspace{3\baselineskip}\vspace{1pt}\textbf{\color{acento}#1}\enspace}

% Fórmula clave, aviso de errores y ejemplo.
\tcbset{before upper=\RaggedRight}
\newtcolorbox{formula}{enhanced,colback=acento!6,colframe=acento!70,boxrule=0.5pt,arc=1.5pt,left=3pt,right=3pt,top=0pt,bottom=0pt,before skip=2pt,after skip=2pt}
\newtcolorbox{ojo}{enhanced,colback=aviso!5,colframe=aviso!60,boxrule=0.5pt,arc=1.5pt,left=3pt,right=3pt,top=1pt,bottom=1pt,before skip=3pt,after skip=2pt,
  fonttitle=\bfseries\footnotesize,coltitle=aviso,colbacktitle=aviso!5,title={Ojo},attach title to upper={\enspace}}
\newtcolorbox{ejemplo}{enhanced,colback=white,colframe=acento!45,boxrule=0.4pt,arc=1.5pt,left=3pt,right=3pt,top=1pt,bottom=1pt,before skip=3pt,after skip=2pt,
  fonttitle=\bfseries\footnotesize,coltitle=acento,colbacktitle=white,title={Ejemplo},attach title to upper={\enspace}}
\newcommand{\dx}{\,dx}

\bloque{algebra}
\begin{document}
\cabecera{5}{Sistemas de ecuaciones}{Lecciones 5.1--5.11: sistemas lineales, soluciones, clasificación, notación matricial, sistemas equivalentes, teorema de Rouché-Frobenius, resolución por la inversa, Gauss y Cramer, sistemas homogéneos y discusión con parámetros.}
\begin{multicols}{2}

\seccion{1. Sistemas lineales y soluciones}
\begin{formula}
\[a_1x_1+a_2x_2+\cdots+a_nx_n=b\]
\end{formula}
Un sistema reúne varias ecuaciones lineales. Una \textbf{solución} es un valor de las incógnitas que cumple \textbf{todas} a la vez. Resolver es encontrar todas las soluciones (o justificar que no hay).
\sub{Comprobar no es resolver:} una solución encontrada no prueba que sea la única.

\seccion{2. Clasificación}
\begin{center}
\small\begin{tabular}{@{}l@{\ \ }l@{\ \ }l@{}}
\textbf{Tipo} & \textbf{Soluciones} & \textbf{Dos rectas}\\ \hline
SCD (compatible determinado) & Una & Secantes\\
SCI (compatible indeterminado) & Infinitas & Coincidentes\\
SI (incompatible) & Ninguna & Paralelas
\end{tabular}
\end{center}
Un sistema lineal nunca tiene exactamente dos soluciones. En un SCI se expresan con parámetros: \(x+y=5\Rightarrow S=\{(t,\,5-t):t\in\mathbb R\}\).

\seccion{3. Notación matricial}
\begin{formula}
\[AX=B,\qquad (A\,|\,B)=\text{matriz ampliada}\]
\end{formula}
\(A\): coeficientes; \(X\): incógnitas; \(B\): términos independientes.
\begin{ejemplo}
\[\begin{cases}x+y+z=6\\2x-y+z=3\\x+2y-z=2\end{cases}\iff \begin{pmatrix}1&1&1\\2&-1&1\\1&2&-1\end{pmatrix}\!\begin{pmatrix}x\\y\\z\end{pmatrix}=\begin{pmatrix}6\\3\\2\end{pmatrix}\]
\end{ejemplo}

\seccion{4. Sistemas equivalentes}
Tienen las mismas soluciones. Transformaciones que lo conservan: intercambiar ecuaciones, multiplicar una ecuación por un número \(\ne0\), sumar a una ecuación un múltiplo de otra. Son la base del método de Gauss.

\seccion{5. Teorema de Rouché-Frobenius}
\begin{formula}
Con \(n\) incógnitas:
\[\operatorname{rg}(A)=\operatorname{rg}(A|B)=n\ \Rightarrow\ \text{SCD}\]
\[\operatorname{rg}(A)=\operatorname{rg}(A|B)<n\ \Rightarrow\ \text{SCI}\]
\[\operatorname{rg}(A)\ne\operatorname{rg}(A|B)\ \Rightarrow\ \text{SI}\]
\end{formula}
En un SCI hay \(n-\operatorname{rg}(A)\) parámetros libres.

\seccion{6. Método de Gauss}
Se escalona la matriz ampliada con transformaciones de filas y se resuelve de abajo arriba.
\begin{ejemplo}
\[\left(\begin{array}{rrr|r}1&1&1&6\\2&-1&1&3\\1&2&-1&2\end{array}\right)\xrightarrow[F_3-F_1]{F_2-2F_1}\left(\begin{array}{rrr|r}1&1&1&6\\0&-3&-1&-9\\0&1&-2&-4\end{array}\right)\]
\[\xrightarrow{3F_3+F_2}\left(\begin{array}{rrr|r}1&1&1&6\\0&-3&-1&-9\\0&0&-7&-21\end{array}\right)\]
\(z=3\), \(-3y-3=-9\Rightarrow y=2\), \(x=6-2-3=1\). SCD: \((1,2,3)\).
\end{ejemplo}
\sub{Al escalonar:} una fila \((0\ 0\ 0\,|\,k)\) con \(k\ne0\) indica SI; una fila de ceros \((0\ 0\ 0\,|\,0)\) se elimina y puede dar un SCI.

\seccion{7. Regla de Cramer}
Para un sistema con tantas ecuaciones como incógnitas y \(\Delta=\det(A)\ne0\) (SCD):
\begin{formula}
\[x_i=\frac{\Delta_i}{\Delta},\qquad \Delta_i=\det(A_i)\]
\end{formula}
\(A_i\) es \(A\) con la columna \(i\) sustituida por \(B\).
\begin{ejemplo}
En el sistema anterior, \(\Delta=7\) y
\[\Delta_x=\begin{vmatrix}6&1&1\\3&-1&1\\2&2&-1\end{vmatrix}=7\Rightarrow x=\frac77=1.\]
\end{ejemplo}

\seccion{8. Resolución con la matriz inversa}
\begin{formula}
\[AX=B\ \Rightarrow\ X=A^{-1}B\qquad(\det A\ne0)\]
\end{formula}
\begin{ejemplo}
\(\begin{cases}2x+y=3\\x+y=2\end{cases}\): \(A^{-1}=\begin{pmatrix}1&-1\\-1&2\end{pmatrix}\), \(X=A^{-1}\begin{pmatrix}3\\2\end{pmatrix}=\begin{pmatrix}1\\1\end{pmatrix}\).
\end{ejemplo}

\seccion{9. Sistemas homogéneos}
Todos los términos independientes son 0: \(AX=O\). Siempre tienen la solución trivial \((0,\dots,0)\).
\begin{formula}
\(\det(A)\ne0\iff\) solo la solución trivial.\\
\(\det(A)=0\iff\) infinitas soluciones (no triviales).
\end{formula}
\begin{ejemplo}
\(\begin{cases}x+2y=0\\2x+4y=0\end{cases}\): \(\det A=0\), \(S=\{(-2t,\,t):t\in\mathbb R\}\).
\end{ejemplo}

\seccion{10. Discusión con parámetros}
\begin{enumerate}
  \item Calcula \(\det(A)\) y halla los valores que lo anulan.
  \item Para el resto de valores: \(\operatorname{rg}(A)=\operatorname{rg}(A|B)=n\), SCD.
  \item Para cada valor crítico, calcula \(\operatorname{rg}(A)\) y \(\operatorname{rg}(A|B)\) y aplica Rouché.
\end{enumerate}
\begin{ejemplo}
\(\begin{cases}x+y=1\\2x+ky=m\end{cases}\): \(\det A=k-2\).
\begin{itemize}
  \item \(k\ne2\): SCD.
  \item \(k=2\), \(m=2\): las ecuaciones son proporcionales, \(\operatorname{rg}(A)=\operatorname{rg}(A|B)=1<2\): SCI.
  \item \(k=2\), \(m\ne2\): \(\operatorname{rg}(A)=1\ne2=\operatorname{rg}(A|B)\): SI.
\end{itemize}
\end{ejemplo}

\seccion{11. Elegir el método}
\begin{tabular}{@{}p{3.6cm}p{4.2cm}@{}}
\textbf{Situación} & \textbf{Método}\\ \hline
Discutir (cuántas soluciones) & Rangos y Rouché\\
Resolver cualquier sistema & Gauss\\
\(n\times n\) con \(\det A\ne0\), pocas incógnitas & Cramer o inversa\\
Varios sistemas con la misma \(A\) & Inversa\\
Parámetro en los coeficientes & \(\det A\), luego casos
\end{tabular}

\seccion{12. Planteamiento de problemas}
Define las incógnitas con unidades, traduce cada condición en una ecuación, resuelve y \textbf{comprueba en el enunciado} que la solución tiene sentido (por ejemplo, cantidades enteras y no negativas).
\begin{ejemplo}
Entradas de adulto (\(x\)) e infantiles (\(y\)): 50 entradas y 380\,euros, a 10\,euros y 6\,euros.
\(x+y=50\), \(10x+6y=380\Rightarrow x=20\), \(y=30\).
\end{ejemplo}

\seccion{13. Tres incógnitas: planos}
Cada ecuación con tres incógnitas es un plano. SCD: los planos se cortan en un punto. SCI: tienen en común una recta (un parámetro) o un plano (dos parámetros). SI: no hay ningún punto común a los tres.

\seccion{14. Discusión 3\texorpdfstring{\(\times\)}{x}3 resuelta}
\[\begin{cases}x+y+z=2\\2x+3y+z=3\\kx+10y+4z=11\end{cases}\]
\begin{enumerate}
  \item \(\det A=\begin{vmatrix}1&1&1\\2&3&1\\k&10&4\end{vmatrix}=14-2k=0\iff k=7\).
  \item \(k\ne7\): \(\operatorname{rg}(A)=\operatorname{rg}(A|B)=3\): SCD.
  \item \(k=7\): \(\begin{vmatrix}1&1\\2&3\end{vmatrix}=1\ne0\Rightarrow\operatorname{rg}(A)=2\). La tercera ecuación es la primera más tres veces la segunda (también en los términos independientes: \(2+9=11\)), así que \(\operatorname{rg}(A|B)=2\): SCI.
  \item Solución para \(k=7\): restando, \(y-z=-1\). Con \(z=t\): \(y=t-1\), \(x=3-2t\).
\end{enumerate}
\[S=\{(3-2t,\ t-1,\ t):t\in\mathbb R\}\]
\begin{ejemplo}
Otro determinante con parámetro: \(\begin{vmatrix}1&1&1\\1&k&1\\1&1&k\end{vmatrix}=(k-1)^2\). Con términos independientes \(1,2,3\): si \(k\ne1\), SCD; si \(k=1\), las dos primeras ecuaciones dicen \(x+y+z=1\) y \(x+y+z=2\): SI.
\end{ejemplo}

\seccion{15. Tres incógnitas en contexto}
Una unidad de cada uno de tres productos cuesta 10 euros; dos de A y una de B, 8 euros; una de A y dos de C, 13 euros. Precios \(x,y,z\) en euros:
\[\begin{cases}x+y+z=10\\2x+y=8\\x+2z=13\end{cases}\]
Restando las dos primeras: \(x-z=-2\). Con la tercera: \(3x=9\Rightarrow x=3\); \(z=5\), \(y=2\). Comprobación: \(3+2+5=10\), \(6+2=8\), \(3+10=13\). A cuesta 3 euros, B 2 euros y C 5 euros.

\seccion{16. Interpretar un SCI}
En \(S=\{(3-2t,\ t-1,\ t)\}\), cada valor de \(t\) da una solución. Si las incógnitas son cantidades, se añaden condiciones: \(3-2t\ge0\), \(t-1\ge0\), \(t\ge0\Rightarrow1\le t\le1{,}5\). Si además deben ser enteras, solo queda \(t=1\): \((1,0,1)\).

\seccion{17. Esquema final}
\begin{tabular}{@{}p{3.3cm}p{4.5cm}@{}}
\textbf{Si\ldots} & \textbf{Entonces\ldots}\\ \hline
\(\det A\ne0\) (\(n\times n\)) & SCD; Cramer o inversa\\
\(\det A=0\) & SCI o SI: comparar rangos\\
\(\operatorname{rg}A=\operatorname{rg}(A|B)<n\) & SCI con \(n-\operatorname{rg}A\) parámetros\\
\(\operatorname{rg}A<\operatorname{rg}(A|B)\) & SI\\
Homogéneo & Siempre compatible
\end{tabular}

\begin{ojo}
\begin{itemize}
  \item Aplicar Cramer o la inversa con \(\det A=0\).
  \item Concluir SI porque \(\det A=0\): también puede ser SCI.
  \item En Gauss, operar solo en \(A\) y olvidar la columna de \(B\).
  \item Dar un ejemplo de solución cuando se pide la familia completa.
  \item Olvidar comprobar la solución en todas las ecuaciones.
\end{itemize}
\end{ojo}

\end{multicols}
\end{document}
